\documentclass[10pt,a4paper]{amsart}
\usepackage[margin=19mm]{geometry}
\usepackage{amsmath,amssymb,mathtools}
\usepackage[scr=rsfs,cal=euler]{mathalfa}
\usepackage{xfrac}
\usepackage[unicode,hidelinks,bookmarks=false]{hyperref}
\newcommand{\ie}{i.e.\ }
\newcommand{\cf}{cf.\ }
\newcommand{\resp}{respectively\ }
\newcommand{\Subsection}{\subsection}
\providecommand{\nobreakdash}{\nobreak\hbox{-}}
\setlength{\emergencystretch}{2em}
\multlinegap0pt
\usepackage{amssymb}
\usepackage{mathtools}
\usepackage{tikz}
\newtheorem{theorem}{Theorem}
\DeclareMathOperator{\argu}{arg}
\newcommand{\bbC}{\mathbb{C}}
\newcommand{\bbR}{\mathbb{R}}
\DeclareMathOperator{\dist}{dist}
\title{New Perspectives on the Unimodality of Domination Polynomials}
\author{Mohamed Omar}
\date{Source: arXiv:2601.14494v1 (CC BY 4.0)}
\begin{document}
\maketitle

\noindent\emph{Source and licence.} New Perspectives on the Unimodality of Domination Polynomials. Version arXiv:2601.14494v1 (CC BY 4.0), distributed by arXiv under the Creative Commons Attribution 4.0 International licence, http://creativecommons.org/licenses/by/4.0/, as recorded in the arXiv licence metadata for this exact version and shown on the abstract page https://arxiv.org/abs/2601.14494v1. Full licence notice: \texttt{licenses/arxiv-2601.14494-CC-BY-4.0.txt}.\par\smallskip

\noindent\textit{Editorial scope.} Counted unit: the article's theorem thm:angularSeparation (source0.tex lines 251--253). The article's complete original proof of it is delivered with it (source0.tex lines 255--276, one \texttt{proof} environment of 22 lines). No mathematics is written, repaired or re-derived: the statement and the proof are the article's own lines, in the article's own notation, and no retained line is substituted or re-flowed.  Retained uncounted as the source-local context the proof needs: lines 237--239.  Nothing was omitted from any retained span, so the package carries no omission record.  No retained line contains a citation, so the wrapper carries no bibliography and no bibliography entry.  No retained line contains a figure environment, an image-inclusion command or drawing code, so no image is delivered and the package declares no asset dependency.  Editorial formatting only, in addition: an AMS \texttt{amsart} preamble that restates the article's own declarations and macros used by the retained lines, namely the environments \texttt{theorem} on separate counters, together with the standard packages those lines need.  The retained environments print in this document as Theorem 1 in order of appearance, whereas the article prints the same items with the numbering of its own full document.  The complete native source of the article as distributed in the arXiv e-print for this exact version is bundled unchanged as \texttt{sources/193090/source0.tex}.
\begin{theorem}\label{thm:rayZeroFree}
Fix an integer $\Delta\ge 2$ and define $T_{\Delta}=\frac{(\Delta-1)^{\Delta+1}}{\Delta^{\Delta}}.$ There exists an open set $\Omega_{\Delta}\subseteq \bbC$ containing the ray $(T_{\Delta},\infty)$ such that for every graph $G$ with $\Delta(G)\le \Delta$ one has $D(G,z)\neq 0$ for all $z\in \Omega_{\Delta}$.
\end{theorem}

\begin{theorem}\label{thm:angularSeparation}
Fix $\Delta\ge 2$ and $r>T_{\Delta}$. There exists $\theta=\theta(\Delta,r)>0$ such that for every graph $G$ with $\Delta(G)\le \Delta$ the domination polynomial $D(G,x)$ has no root $z$ satisfying $|z|\ge r$ and $|\argu(z)|<\theta$.
\end{theorem}

\begin{proof}
Let $M=\Delta+1$ and $U_M$ be the open set from the proof of Theorem~\ref{thm:rayZeroFree}. Since $r>T_{\Delta}$, we have $1/r<L_{\Delta}$, where $L_{\Delta}=1/T_{\Delta}$. Consider the compact interval
\[
K:=\left[0,\frac{1}{r}\right]\subseteq \bbR.
\]
Because $U_M$ is open and contains $K$, there exists $\varepsilon>0$ such that the open $\varepsilon$-neighborhood of $K$ is contained in $U_M$. In other words, if $\lambda\in\bbC$ satisfies $\dist(\lambda,K)<\varepsilon$, then $\lambda\in U_M$. Now choose $\theta>0$ such that
\begin{equation}\label{eq:thetaChoice}
\frac{2}{r}\sin(\theta/2)<\varepsilon.
\end{equation}
We show that for any $z\in\bbC$ satisfying $|z|\ge r$ and $|\argu(z)|<\theta$, $D(G,z) \neq 0$. Write $z=|z|e^{i\phi}$ with $|\phi|<\theta$. Then $\frac{1}{z}=\frac{1}{|z|}e^{-i\phi}.$ Let $t:=1/|z|$, so $0<t\le 1/r$ and hence $t\in K$. We estimate
\[
\left|\frac{1}{z}-t\right|
=t\left|e^{-i\phi}-1\right|
=t\cdot 2\sin(|\phi|/2)
\le \frac{2}{|z|}\sin(\theta/2)
\le \frac{2}{r}\sin(\theta/2)
<\varepsilon
\]
by \eqref{eq:thetaChoice}. Thus $\dist(1/z,K)<\varepsilon$, so $1/z\in U_M$. Equivalently $z\in\Omega_{\Delta}$, where $\Omega_{\Delta}$ is as in Theorem~\ref{thm:rayZeroFree}.

Now let $G$ satisfy $\Delta(G)\le \Delta$. By Theorem~\ref{thm:rayZeroFree}, $D(G,z)\neq 0$ for all $z\in\Omega_{\Delta}$. In particular, $D(G,z)\neq 0$ for all $z$ with $|z|\ge r$ and $|\argu(z)|<\theta$.
\end{proof}

\end{document}
